\documentclass[UTF8,12pt,a4paper]{ctexart}
\usepackage[a4paper,margin=2.15cm]{geometry}
\usepackage{amsmath,amssymb,bm,booktabs,array,longtable,multirow}
\usepackage{enumitem}
\usepackage{fancyhdr}
\usepackage{lastpage}
\linespread{1.3}
\setlength{\parindent}{2em}
\setlength{\parskip}{0.2em}
\setlist[enumerate]{leftmargin=2.5em,itemsep=0.45em,topsep=0.3em}
\setlist[itemize]{leftmargin=2.2em,itemsep=0.25em,topsep=0.2em}
\renewcommand\arraystretch{1.18}
\newcommand{\blank}[1][2.4cm]{\underline{\hspace{#1}}}
\newcommand{\dd}{\mathrm{d}}
\newcommand{\ee}{\mathrm{e}}
\newcommand{\jj}{\mathrm{j}}
\newcommand{\prob}{\mathsf{P}}
\pagestyle{fancy}
\fancyhf{}
\cfoot{第\thepage 页 / 共\pageref{LastPage}页}
\renewcommand{\headrulewidth}{0pt}
\begin{document}

\begin{center}
    {\zihao{2}\bfseries 误差理论与数据处理试卷（三）}\par
    \vspace{0.8em}
\end{center}
\vspace{0.6em}

\section*{一、填空题答案}
\begin{enumerate}[label=\arabic*、]
\item 单峰性、有界性、对称性、抵偿性；抵偿性。
\item $4.510$，$6.379$。
\item $-0.002$，实验对比，$0.002$，$10.004$。
\item 绝对，相对。
\item $50\mu\Omega$，B，$8$。
\item 
\begin{equation*}
A=\begin{bmatrix}1&-3\\4&-1\\2&-1\end{bmatrix},\quad
L=\begin{bmatrix}-5.6\\8.1\\0.5\end{bmatrix},\quad
P=\begin{bmatrix}1&0&0\\0&2&0\\0&0&3\end{bmatrix}.
\end{equation*}
\item $31.5$。
\end{enumerate}

\section*{二、选择题答案}
\begin{center}
\begin{tabular}{ccccccccc}
\toprule
1&2&3&4&5&6&7&8&9\\
\midrule
A&B、F&C&C&A&B&C&D&D\\
\bottomrule
\end{tabular}
\end{center}

\section*{三、判断题答案}
\begin{center}
\begin{tabular}{cccccccccc}
\toprule
1&2&3&4&5&6&7&8&9&10\\
\midrule
错&对&错&错&错&对&对&错&错&错\\
\bottomrule
\end{tabular}
\end{center}

\newpage

\section*{四、计算题答案}
\begin{enumerate}[label=\arabic*、]
\item 记频数为 $p_i$，则
\begin{equation*}
    \bar{R}=\frac{\sum p_iR_i}{\sum p_i}=1215.06\Omega.
\end{equation*}
频数总和 $N=200$，单次测量标准差为
\begin{equation*}
    s=\sqrt{\frac{\sum p_i(R_i-\bar{R})^2}{N-1}}
     =\sqrt{\frac{509.28}{199}}=1.600\Omega.
\end{equation*}
平均值的标准差为
\begin{equation*}
    s_{\bar{R}}=s/\sqrt{N}=1.600/\sqrt{200}=0.113\Omega.
\end{equation*}

\item 放大率最可信赖值为
\begin{equation*}
    D=f_1/f_2=19.80/0.800=24.75.
\end{equation*}
灵敏系数为
\begin{equation*}
    c_1=\frac{\partial D}{\partial f_1}=\frac{1}{f_2}=1.25,
    \qquad
    c_2=\frac{\partial D}{\partial f_2}=-\frac{f_1}{f_2^2}=-30.94.
\end{equation*}
由 $f_1$ 引起的不确定度分量为
\begin{equation*}
    u_1=|c_1|\sigma_1=1.25\times0.10=0.125.
\end{equation*}
由 $f_2$ 引起的不确定度分量为
\begin{equation*}
    u_2=|c_2|\sigma_2=30.94\times0.005=0.155.
\end{equation*}
放大率 $D$ 的标准不确定度为
\begin{equation*}
    u_c=\sqrt{u_1^2+u_2^2}=\sqrt{0.125^2+0.155^2}=0.199.
\end{equation*}
有效自由度按 Welch--Satterthwaite 公式估计为
\begin{equation*}
    \nu_{\mathrm{eff}}=\frac{u_c^4}{u_1^4/8+u_2^4/4}\approx9.
\end{equation*}
当 $P=95\%$ 时，取 $t_{0.05}(9)=2.26$，故
\begin{equation*}
    U_{95}=2.26\times0.199=0.45.
\end{equation*}
所以可写为 $D=24.75\pm0.45$，$P=95\%$。

\newpage

\item 设
\begin{equation*}
A=\begin{bmatrix}2&1\\1&-1\\4&-1\\1&4\end{bmatrix},
\qquad
L=\begin{bmatrix}5.1\\1.1\\7.4\\5.9\end{bmatrix}.
\end{equation*}
则
\begin{equation*}
A^TA=\begin{bmatrix}22&1\\1&19\end{bmatrix},
\qquad
(A^TA)^{-1}=\frac{1}{417}\begin{bmatrix}19&-1\\-1&22\end{bmatrix}.
\end{equation*}
最小二乘估计为
\begin{equation*}
    \hat{X}=(A^TA)^{-1}A^TL=\begin{bmatrix}2.08\\0.95\end{bmatrix}.
\end{equation*}
按原题答案取两位小数计算残差：
\begin{equation*}
\begin{aligned}
    v_1&=5.1-(2\times2.08+0.95)=-0.01,\\
    v_2&=1.1-(2.08-0.95)=-0.03,\\
    v_3&=7.4-(4\times2.08-0.95)=0.03,\\
    v_4&=5.9-(2.08+4\times0.95)=0.02.
\end{aligned}
\end{equation*}
故
\begin{equation*}
    \sum v_i^2=23\times10^{-4},
    \qquad
    \sigma=\sqrt{\frac{23\times10^{-4}}{4-2}}=0.034.
\end{equation*}
由 $d_{11}=19/417=0.046,d_{22}=22/417=0.053$，得
\begin{equation*}
    \sigma_x=0.034\sqrt{0.046}=0.007,
    \qquad
    \sigma_y=0.034\sqrt{0.053}=0.008.
\end{equation*}

\newpage

\item 设经验公式为 $\hat{y}=b_0+bx$。由数据得
\begin{equation*}
\sum x_i=180,\quad \sum y_i=692,\quad \sum x_i^2=3450,\quad
\sum y_i^2=49470,\quad \sum x_iy_i=13032.
\end{equation*}
因此
\begin{equation*}
\bar{x}=18,\quad \bar{y}=69.2,
\end{equation*}
\begin{equation*}
l_{xx}=\sum x_i^2-\frac{1}{n}\left(\sum x_i\right)^2=210,
\quad
l_{yy}=\sum y_i^2-\frac{1}{n}\left(\sum y_i\right)^2=1583.6,
\end{equation*}
\begin{equation*}
l_{xy}=\sum x_iy_i-\frac{1}{n}\left(\sum x_i\right)\left(\sum y_i\right)=576.0.
\end{equation*}
回归系数为
\begin{equation*}
    b=l_{xy}/l_{xx}=2.74,
    \qquad b_0=\bar{y}-b\bar{x}=19.8.
\end{equation*}
故经验公式为
\begin{equation*}
    \hat{y}=19.8+2.74x.
\end{equation*}
方差分析：
\begin{equation*}
    U=bl_{xy}=1578.24,
    \qquad Q=l_{yy}-U=5.36,
    \qquad S=l_{yy}=1583.6.
\end{equation*}
\begin{center}
\begin{tabular}{ccccc}
\toprule
来源 & 平方和 & 自由度 & 方差 & $F$ \\
\midrule
回归 & $U=1578.24$ & $1$ & -- & -- \\
残余 & $Q=5.36$ & $8$ & $Q/8=0.67$ & $F=1578.24/0.67=2356$ \\
总计 & $S=1583.6$ & $9$ & -- & -- \\
\bottomrule
\end{tabular}
\end{center}
因为
\begin{equation*}
    F=2356>F_{0.01}(1,8)=11.26,
\end{equation*}
所以回归方程在 $0.01$ 水平上高度显著。

\end{enumerate}
\end{document}
